% !TeX root = ../main.tex
\section{Conclusions}
\label{sec:conclusion}

\textcolor{red}{This paper presented a room-level conditional graph framework for shear wall layout generation in high-rise residential buildings.} By treating rooms and room adjacencies as the basic units of prediction, the method aligns graph representation with architectural space semantics and with the boundary-constrained nature of shear wall placement. FiLM-based conditioning and dual-stream training further allow the model to respond to prescribed density regimes under sparse paired supervision.

Experiments on 143 high-rise residential floor plans show that this formulation achieves a fold-averaged Image IoU of 0.565 $\pm$ 0.005 and an ensemble Image IoU of 0.597, exceeding the component-based GNN baseline in both accuracy and stability. Within this dataset, the main empirical finding is not only that the room-level representation improves overall accuracy, but also that it reduces spurious wall predictions and strengthens condition responsiveness across the three tested discrete density groups. Finite-element case studies on two real projects provide additional, case-based evidence that the generated schemes can preserve acceptable global stiffness characteristics under the tested settings.

The current framework still has several limitations, corresponding to the issues discussed above. Its validation is limited to regular high-rise residential layouts, so broader typological transfer requires further testing, especially for plans with highly irregular regions or strong segment-level constraints. The method also depends on reliable CAD preprocessing and does not directly optimize structural response indicators during training. In addition, the condition input is represented by discrete density groups, and the present validation assumes a repeated standard-floor layout rather than explicitly modeling vertical layout variation.

Future work should focus on three directions that follow directly from the current limits of the framework: supporting irregular polygonal rooms without decomposition, replacing discrete density categories with continuous engineering descriptors, and coupling generation more tightly with structural-performance objectives during training.
